\documentclass[12pt]{article}

\usepackage{geometry}
\geometry{margin = 1in,top =0.12\paperheight, headheight=\paperheight}
\usepackage[export]{adjustbox}
\usepackage{array}
\usepackage{amsmath}
\usepackage{amsfonts}
\usepackage{fancyhdr}
\usepackage{lastpage}
\usepackage{xcolor}
\usepackage{comment}
%\renewcommand{\footrulewidth}{0.4pt}

\usepackage{enumitem}
\usepackage{pifont}
\usepackage{graphicx}
\graphicspath{{../img/}}

\newtheorem{theorem}{Theorem}
\newtheorem{exercise}{Exercise}


\newcommand{\R}{\mathbb R}
\newcommand{\e}{{\rm e}}
\newcommand{\inpr}[1]{\left\langle#1\right\rangle}
\newcommand{\norm}[1]{\lVert #1 \rVert}
\newcommand{\abs}[1]{\lvert #1 \rvert}
\newcommand{\vv}{\mathbf v}
\newcommand{\uv}{\mathbf u}

\DeclareMathOperator{\xd}{d\!}
\DeclareMathOperator{\proj}{proj}
\pagestyle{fancy}
\fancyhf{}
\rhead{Written Assignment, Page \thepage}
\lhead{MATH211}
\chead{\includegraphics[width = 0.15\textwidth]{MCLogo-Bck.png}}

\begin{document}
\noindent
{\bf Problem.}
Since $\mathbf a = a_{\mathbf T}\mathbf T + a_{\mathbf N}\mathbf N$, $a_{\mathbf N}\mathbf N = \mathbf a - a_{\mathbf T}\mathbf T$. In the class, we proved that $a_{\mathbf N} = v^2K > 0$, thus $a_{\mathbf N}\mathbf N$ and $\mathbf N$ are in the same direction. Based on these observations, continue to prove that
\[
	\mathbf N = \frac{\left(\mathbf v\cdot\mathbf v\right)\mathbf a - \left(\mathbf v\cdot\mathbf a\right)\mathbf v}{\norm{\left(\mathbf v\cdot\mathbf v\right)\mathbf a - \left(\mathbf v\cdot\mathbf a\right)\mathbf v}}.
\]
You can use any theorem that is in the textbook or proved in the class.
\end{document}